\section{Data Modeling}
\begin{justify}
The application uses three main entities: users, tasks, and task analysis. A user can have many tasks. Each task can have one analysis record containing difference, accuracy, estimation type, and analysis timestamp.
\end{justify}

\begin{table}[H]
\centering
\begin{tabular}{|M|L|}
\hline
\textbf{Table} & \textbf{Fields} \\
\hline
users & user\_id, name, email, password, created\_at \\
\hline
tasks & task\_id, user\_id, task\_name, category, expected\_time, actual\_time, status, created\_at \\
\hline
task\_analysis & analysis\_id, task\_id, difference, accuracy, estimation\_type, analyzed\_at \\
\hline
\end{tabular}
\caption{Data Dictionary}
\end{table}

\section{Activity Diagrams / Class Diagram}
\begin{justify}
The activity flow begins when the user opens the frontend and fills the task form. The frontend sends the data to the backend API. The backend validates the request, performs rule-based analysis, inserts task data into the database, inserts analysis data, updates task status, and returns the result to the frontend.
\end{justify}

\section{Functional Modeling}
\begin{justify}
The Level-0 DFD contains one external entity, User, and one main process, AI-Based Task Estimation System. The User provides task details and receives analysis results. The system communicates with the database to store and retrieve user, task, and analysis data.
\end{justify}

\section{Architectural Design}
\begin{justify}
The project follows a three-layer architecture. The presentation layer contains HTML, CSS, and JavaScript. The application layer contains Express.js APIs and rule-based AI logic. The data layer contains MySQL tables and SQL queries.
\end{justify}
